\documentclass[11pt,paper=a4,answers]{exam}
\usepackage{graphicx,lastpage}
\usepackage{upgreek}
\usepackage{censor}
\usepackage{gensymb}
\usepackage{amsmath}
\usepackage{multicol}
\usepackage{enumitem}
\usepackage{tasks}
\usepackage{adjustbox}
\usepackage{xcolor}
\usepackage{tfrupee}
\setlength{\voffset}{0.25in}
\usepackage{tabularx}
\newcolumntype{Y}{>{\centering\arraybackslash}X}
%==============================================================
% Customise Non-Essential Packages Below
% =============================================================
\usepackage[most]{tcolorbox}
\usepackage{eso-pic}
\usepackage{tikz}
\AddToShipoutPictureBG{%
\begin{tikzpicture}[remember picture,overlay]
\foreach \x in {-8,-4,0,4,8,12,16,20}
\foreach \y in {-8,-4,0,4,8,12,16,20} {
\node[opacity=0.05,rotate=45,anchor=center] at ([xshift=\x cm,yshift=\y cm]current page.center) {%
\begin{minipage}{2cm}
\centering
\includegraphics[width=2cm]{images/logo_black.png}\\
\small preppeo.com
\end{minipage}%
};
}
\end{tikzpicture}%
}
\printanswers

%==============================================================
% Customise Non-Essential Packages Above
% =============================================================
\definecolor{svgray}{rgb}{0.90, 0.90, 0.90}
\graphicspath{{images/}}

\usepackage[normalem]{ulem}
\renewcommand\ULthickness{1pt}   %%---> For changing thickness of underline
\setlength\ULdepth{3.5ex}%\maxdimen ---> For changing depth of underline
\renewcommand{\baselinestretch}{1}
\chead{
\includegraphics[scale=0.1,height=2cm ]{logo.png}
}
\pagestyle{empty}

\pagestyle{headandfoot}
\headrule
\cfoot{W: https://preppeo.com; M: +91-8130711689}

\begin{document}
%==============================================================
% Customise Question paper details below this
% =============================================================
\begin{center}
    \centering
    \renewcommand{\arraystretch}{1.5}
    \begin{tabularx}{\textwidth}{YYY}
    & \normalsize\bf{{\Large{{Quant}}}} & \\
    By Preppeo & \normalsize\bf{{geo\_sat\_hard\_spr\_006}} & SAT\\
    \normalsize{{\textbf{{Algebra}} }} & \normalsize{{\textbf{{Area volumes, circle, triangle }} }} & \normalsize{{\textbf{{Solving Linear Equations}} }} \\
    \end{tabularx}
\end{center}
\par\noindent\rule{\textwidth}{1pt}
% =============================================================
% Customise Question paper details above this
% =============================================================

\begin{questions}
\pointsinrightmargin
\pointsdroppedatright
\marksnotpoints
\pointpoints{mark}{marks}
\pointformat{\boldmath\themarginpoints}
\bracketedpoints

% =============================================================
% Question Begin
% =============================================================

\section*{SAT Geometry \& Trigonometry - Practice Set 8 (Hard Numeric Entry)}

% Q1: Based on Similar Rectangle Area logic
\question
Rectangles $PQRS$ and $WXYZ$ are similar. The length of each side of $WXYZ$ is $5$ times the length of the corresponding side of $PQRS$. The area of $PQRS$ is $42$ square units. What is the area, in square units, of $WXYZ$?
\begin{center}
\begin{tikzpicture}[scale=0.6]
    \draw[thick] (0,0) rectangle (2,1.2);
    \node at (1,0.6) {\small Area 42};
    \node[below] at (1,0) {\tiny $PQRS$};
    
    \draw[thick] (4,0) rectangle (9,3);
    \node at (6.5,1.5) {\small Area ?};
    \node[below] at (6.5,0) {\tiny $WXYZ$};
    
    \draw[<->] (2.2,0.6) -- (3.8,0.6) node[midway, above]{}; 
\end{tikzpicture}
\end{center}

\begin{solution}
When two figures are similar, the ratio of their areas is the square of the scale factor ($k$).
\vspace{5pt}
\begin{center}
Scale Factor ($k$) $= 5$
\vspace{5pt}
Area Ratio $= k^2 = 5^2 = 25$
\vspace{5pt}
Area of $WXYZ = 42 \times 25$
\vspace{5pt}
$42 \times 25 = 1,050$
\end{center}
\vspace{5pt}
Answer: \colorbox{yellow!100}{1050}
\end{solution}

\vspace{1cm}

% Q2: Based on Percentage Increase Area logic
\question
A rectangular photograph has an area of $200$ square inches. A digital copy of the photograph is made in which the length and width of the original photograph are each increased by $30\%$. What is the area of the digital copy, in square inches?
\begin{center}
\begin{tikzpicture}[scale=0.8]
    \draw[thick, fill=gray!10] (0,0) rectangle (2,1.5);
    \node at (1,0.75) {\tiny $200$ in$^2$};
    
    \draw[thick, dashed] (4,-0.225) rectangle (6.6,1.725);
    \node at (5.3,0.75) {\tiny Area ?};
    \node[below] at (5.3,-0.225) {\tiny +30\% Dimensions};
\end{tikzpicture}
\end{center}

\begin{solution}
If the dimensions are each increased by $30\%$, the new dimensions are $1.3$ times the original.
\vspace{5pt}
\begin{center}
New Area $= \text{Original Area} \times (\text{Scale Factor})^2$
\vspace{5pt}
New Area $= 200 \times (1.3)^2$
\vspace{5pt}
New Area $= 200 \times 1.69 = 338$
\end{center}
\vspace{5pt}
Answer: \colorbox{yellow!100}{338}
\end{solution}

\vspace{1cm}

% Q3: Based on Arc Fraction logic
\question
Points $X$ and $Y$ lie on a circle with radius $3$. Arc $XY$ has a length of $\frac{3\pi}{4}$. What fraction of the circumference of the circle is the length of arc $XY$?
\begin{center}
\begin{tikzpicture}[scale=1.2]
    \draw[thick] (0,0) circle (1);
    \draw[ultra thick, blue] (0:1) arc (0:45:1);
    \draw[dashed] (0,0) -- (0:1);
    \draw[dashed] (0,0) -- (45:1);
    \node[right] at (0:1) {\tiny $X$};
    \node[above right] at (45:1) {\tiny $Y$};
    \node at (0.3,0.1) {\tiny $r=3$};
\end{tikzpicture}
\end{center}

\begin{solution}
Divide the length of the arc by the total circumference ($C = 2\pi r$).
\vspace{5pt}
\begin{center}
Circumference ($C$) $= 2\pi(3) = 6\pi$
\vspace{5pt}
Fraction $= \dfrac{\text{Arc Length}}{\text{Circumference}} = \dfrac{\frac{3\pi}{4}}{6\pi}$
\vspace{5pt}
Fraction $= \dfrac{3\pi}{4} \times \dfrac{1}{6\pi} = \dfrac{3}{24} = \dfrac{1}{8}$
\end{center}
\vspace{5pt}
Answer: \colorbox{yellow!100}{1/8}
\end{solution}

\vspace{1cm}

% Q4: Based on Trigo Similarity (tan A) logic
\question
Triangle $LMN$ is similar to triangle $PQR$, where $L$ corresponds to $P$ and $N$ corresponds to $R$. Angles $N$ and $R$ are right angles. If $\tan(L) = \frac{4}{3}$ and $PR = 60$, what is the length of $\overline{PQ}$?

\begin{solution}
Since the triangles are similar, $\tan(P) = \tan(L) = \frac{4}{3}$.
\vspace{5pt}
\begin{center}
$\tan(P) = \dfrac{\text{Opposite}}{\text{Adjacent}} = \dfrac{QR}{PR}$
\vspace{5pt}
$\dfrac{4}{3} = \dfrac{QR}{60} \implies QR = 80$
\vspace{10pt}
Use the Pythagorean theorem to find the hypotenuse $PQ$:
\vspace{5pt}
$PQ = \sqrt{60^2 + 80^2} = \sqrt{3600 + 6400} = 100$
\end{center}
\vspace{5pt}
Answer: \colorbox{yellow!100}{100}
\end{solution}

\vspace{1cm}

% Q5: Based on Rectangle Diagonal logic
\question
The length of a rectangle's diagonal is $2\sqrt{29}$, and the length of the rectangle's shorter side is $4$. What is the length of the rectangle's longer side?

\begin{solution}
The sides and diagonal form a right triangle. Use $a^2 + b^2 = c^2$.
\vspace{5pt}
\begin{center}
$4^2 + b^2 = (2\sqrt{29})^2$
\vspace{5pt}
$16 + b^2 = 4 \times 29 = 116$
\vspace{5pt}
$b^2 = 116 - 16 = 100$
\vspace{5pt}
$b = \sqrt{100} = 10$
\end{center}
\vspace{5pt}
Answer: \colorbox{yellow!100}{10}
\end{solution}

\vspace{1cm}

% Q6: Based on Trig Scale logic
\question
In the figure below, $\cos(B) = 0.8$ and $BC = 25$. What is the length of $\overline{BD}$?
\begin{center}
\begin{tikzpicture}[scale=0.8]
    \draw[thick] (0,0) coordinate (A) -- (4,0) coordinate (C) -- (0,3) coordinate (B) -- cycle;
    \draw[thick] (0,0) -- (1.44,1.92) coordinate (D);
    \draw (0,0) rectangle (0.2,0.2);
    \node[below left] at (A) {$A$};
    \node[right] at (C) {$C$};
    \node[above] at (B) {$B$};
    \node[above right] at (D) {$D$};
    \node[left] at (0, 1.5) {$25$};
    \node at (0.3, 2.5) {\tiny $B$};
\end{tikzpicture}
\end{center}

\begin{solution}
In right triangle $BDC$, side $BD$ is adjacent to angle $B$, and $BC$ is the hypotenuse.
\vspace{5pt}
\begin{center}
$\cos(B) = \dfrac{\text{Adjacent}}{\text{Hypotenuse}} = \dfrac{BD}{BC}$
\vspace{5pt}
$0.8 = \dfrac{BD}{25} \implies BD = 20$
\end{center}
\vspace{5pt}
Answer: \colorbox{yellow!100}{20}
\end{solution}


% =============================================================
% Question End
% =============================================================

\end{questions}
\begin{center}
\rule{.5\textwidth}{1pt}
\end{center}
\end{document}
